\SourceLeaf{23}
\[
\alpha\delta\xi=\delta\alpha.\int\!\int\frac{\mathrm{d}p.\mathrm{d}q.a\sin p}{r^3}\bigl(\frac{(a-x)x}{\alpha\alpha}+\frac{(b-y)y}{\beta\beta}+\frac{(c-z)z}{\gamma\gamma}\bigr)
\]
Huius aequationis pars ad dextram per aequ. $[3]$ fit vel $=0$ vel $=-\frac{4\pi a\delta\alpha}{\alpha\beta\gamma}$, prout $M$ iacet extra vel intra corpus, ita ut fiat in casu priori
\[
[4]\qquad\delta\xi=0
\]
in posteriori autem
\[
[5]\qquad\delta\xi=-\frac{4\pi a\delta\alpha}{\alpha\alpha\beta\gamma}
\]
\par
Aequatio $[4]$ protinus ostendit, $\xi$ esse constantem, sive attractionem $X$ massae proportionalem pro omnibus ellipsoidibus, in quibus $\alpha\alpha-\beta\beta$, $\alpha\alpha-\gamma\gamma$ sint quantitates constantes, i. e., quarum tres sectiones principales sint ellipses ex iisdem focis descriptae, quamdiu punctum attractum extra sphaeroidem iaceat. Quam conclusionem, quum omni rigore vera sit, quantumvis proxime sphaeroidis superficies ad punctum attractum accedat, necessario etiam ad sphaeroidem ipsam extendere licebit, cuius superficies per ipsum punctum attractum transit.

Problema itaque de attractione sphaeroidis in punctum quodcunque externum determinanda, reducitur ad duo alia problemata, scilicet primo ad determinationem dimensionum alius sphaeroidis ex iisdem quibus sphaerois proposita focis descriptae punctumque attractum transeuntis, secundo ad problema de attractione sphaeroidis in punctum in ipsius superficie positum. Problema prius pendet a solutione aequationis cubicae, quam semper radicem realem unicam involvere facile demonstratur, cuique hic immorari superfluum videtur. Ut vero problema alterum solvamus, consideremus casum alterum, ubi punctum attractum iacet intra corpus. Quum sit $\beta\beta=\alpha\alpha+BB-AA$, $\gamma\gamma=\alpha\alpha+CC-AA$, substituemus hos valores in aequatione $5$, simulque faciemus $\frac{A}{\alpha}=t$. Hinc emergit
\[
\delta\xi=\frac{4a\pi tt\delta t}{A^3\BareRoot{((1-(1-\frac{BB}{AA})tt)(1-(1-\frac{CC}{AA})tt)}}
\]
sive restituendo characteristicam $\mathrm{d}$, et integrando
\[
\xi=\frac{4a\pi}{A^3}\int\frac{tt\mathrm{d}t}{\BareRoot{((1-(1-\frac{BB}{AA})tt)(1-(1-\frac{CC}{AA})tt)}}
\]
quod integrale ita sumendum est, ut evanescat pro $t=0$, ac dein, pro sphaeroide determinata, cuius semiaxes sunt $A,B,C$, extendendum usque ad $t=1$. Habemus itaque
